% !TeX root = ../Harmonia-EA-Capability-Map.tex
% Copyright (c) 2026 Mark Hunter. SPDX-License-Identifier: GPL-3.0-or-later
% Generated from the explicit 0.2-draft model; all assertion status is retained.
\chapter{B2.1 — Individual Care Delivery ↔ Health Information Access}\label{rad:b2-1}
\noindent\textbf{RAD source:} \texttt{\seqsplit{docs/markdown/05-information-systems-architecture/traceability/b2-1-individual-care-delivery-health-information-access.md}}

The following is a presentation transcription of the governed source. Source catalogue layout and retained illustrations do not allocate tiers or establish additional identities. Equations wrap as text; diagram notation is textual.

\phantomsection\label{rad:b2-1:1}
\phantomsection\label{rad:b2-1:3}
\section*{B2.1 — Individual Care Delivery ↔ Health Information Access}
\phantomsection\label{rad:b2-1:5}

\textbf{B2.1 — Individual Care Delivery ↔ Health Information Access: PASS}

\phantomsection\label{rad:b2-1:7}
\subsection*{1. Authority and bounded result}
\phantomsection\label{rad:b2-1:9}

This record reconciles the explicitly agreed Domain05 Step B B2.1 decisions, 2026-10-13. The \hyperref[rad:business:148]{Domain02 Business Capability catalogue} defines Individual Care Delivery and its six L2 capabilities. The established \hyperref[rad:ise:72]{Health Information Access responsibility} defines the enabling anchor. The \href{https://github.com/fhirfactory/harmonia/blob/98410db60155198a75f74c79d98060070209faa1/docs/markdown/05-information-systems-architecture/traceability/step-b-method.md\#capability-enablement-semantics}{Step B method} governs relationship semantics and Health Information Access scope.

\phantomsection\label{rad:b2-1:16}

PASS applies only to the L1 relationship and six L2 relationships below. It is not a runtime verification, a declaration of exhaustive capability coverage or an application/\allowbreak{}component allocation. This record establishes no additional capability decomposition, Feature, Service, Function or Information mapping. \textbf{B2.2 is not commenced.} No other Business Capability or Business Enabling Capability is mapped by this reconciliation.

\phantomsection\label{rad:b2-1:23}
\subsection*{2. L1 relationship and rationale}
\phantomsection\label{rad:b2-1:25}

\texttt{\seqsplit{Individual Care Delivery enabled-by Health Information Access}}

\phantomsection\label{rad:b2-1:27}

Individual Care Delivery requires governed access to relevant longitudinal and current health information to support clinical assessment, diagnosis, treatment decisions, interventions and monitoring. Health Information Access enables access to that information but does not determine its clinical significance or perform the clinical decision-making of Individual Care Delivery.

\phantomsection\label{rad:b2-1:33}
\subsection*{3. L2 relationships and rationales}
\phantomsection\label{rad:b2-1:35}

These are material capability-enablement relationships. They do not transfer clinical accountability or require each runtime exercise to retrieve information through Health Information Access. The L2 labels identify the agreed Domain02 decomposition; no additional structural Canonical ID scheme is established.

\phantomsection\label{rad:b2-1:40}
\subsubsection*{01.01 Individual Health Assessment}
\phantomsection\label{rad:b2-1:42}

\texttt{\seqsplit{Individual Health Assessment enabled-by Health Information Access}}

\phantomsection\label{rad:b2-1:44}

Individual Health Assessment requires governed access to relevant existing and current health information so that available clinical, health and contextual information can be considered when establishing an understanding of the individual's health state. Health Information Access enables retrieval of that information but does not perform the assessment or establish the assessed findings.

\phantomsection\label{rad:b2-1:50}

Health Information Access does not own every mechanism by which assessment information is acquired. Direct examination, observation, history-taking or measurement may establish assessment information independently.

\phantomsection\label{rad:b2-1:54}
\subsubsection*{01.02 Clinical Determination}
\phantomsection\label{rad:b2-1:56}

\texttt{\seqsplit{Clinical Determination enabled-by Health Information Access}}

\phantomsection\label{rad:b2-1:58}

Clinical Determination requires governed access to relevant current and historical health information where that information is necessary to interpret assessed findings and establish clinically meaningful conclusions about the individual's health state, condition, risk or prognosis. Health Information Access enables retrieval of that information but does not interpret its clinical significance or establish the clinical determination.

\phantomsection\label{rad:b2-1:65}
\subsubsection*{01.03 Individual Care Planning}
\phantomsection\label{rad:b2-1:67}

\texttt{\seqsplit{Individual Care Planning enabled-by Health Information Access}}

\phantomsection\label{rad:b2-1:69}

Individual Care Planning requires governed access to relevant current and historical health and care information where that information is necessary to establish an appropriate intended clinical response, including care objectives, interventions and monitoring. Health Information Access enables retrieval of that information but does not establish the care intent, select the clinical response or perform the planning.

\phantomsection\label{rad:b2-1:76}

Information concerning operational availability of services, resources, staff, facilities or appointments is not Health Information Access merely because it influences care planning.

\phantomsection\label{rad:b2-1:80}
\subsubsection*{01.04 Clinical Intervention}
\phantomsection\label{rad:b2-1:82}

\texttt{\seqsplit{Clinical Intervention enabled-by Health Information Access}}

\phantomsection\label{rad:b2-1:84}

Clinical Intervention requires governed access to relevant health and care information where that information is necessary to safely and appropriately perform an intended clinical action, including information concerning the individual's current health state, established care intent and clinically relevant history. Health Information Access enables retrieval of that information but does not authorise, select or perform the clinical intervention.

\phantomsection\label{rad:b2-1:91}

Access to information describing an intended intervention does not confer authority to perform that intervention.

\phantomsection\label{rad:b2-1:94}
\subsubsection*{01.05 Individual Health Monitoring}
\phantomsection\label{rad:b2-1:96}

\texttt{\seqsplit{Individual Health Monitoring enabled-by Health Information Access}}

\phantomsection\label{rad:b2-1:98}

Individual Health Monitoring requires governed access to relevant current and historical health information where that information is necessary to establish change, stability or deviation in an individual's health state over time. Health Information Access enables access to prior, baseline and current health information required for comparison, but does not perform the monitoring or interpret the clinical significance of an identified change.

\phantomsection\label{rad:b2-1:105}
\subsubsection*{01.06 Individual Care Review}
\phantomsection\label{rad:b2-1:107}

\texttt{\seqsplit{Individual Care Review enabled-by Health Information Access}}

\phantomsection\label{rad:b2-1:109}

Individual Care Review requires governed access to relevant health and care information in order to evaluate delivered care and observed health progression against the intended care objectives. Health Information Access enables access to the care intent, relevant care history and health outcomes required for that evaluation, but does not perform the review or establish the continuing appropriateness of the care intent.

\phantomsection\label{rad:b2-1:116}
\subsection*{4. Preserved boundaries}
\phantomsection\label{rad:b2-1:118}

The \hyperref[rad:business:71]{Domain02 decomposition and context rules} apply: the six established children are now explicitly reconciled as BC CT2s. The current catalogue establishes the stated assessment, monitoring and review CT3s; the other branches retain their recorded uncertainty/\allowbreak{}stopping points. Under the \hyperref[rad:metamodel:33]{fixed Capability model}, subsequently required Harmonia Features must first have meaningful \textbf{Business Enabling CT3} frames established. BC decomposition does not terminate in Features and this B2.1 support trace establishes no ancestry or new CT3 enablement edge. None may be invented for structural completeness. Retained L1/\allowbreak{}L2 display labels alone do not allocate CT tiers. Examples are not subordinate Capabilities, and care setting, delivery mode, temporal pattern and organisational arrangement do not themselves define different business abilities. Acute, Planned/\allowbreak{}Scheduled and Longitudinal Care are not the ICD L2 decomposition.

\phantomsection\label{rad:b2-1:133}

\href{https://github.com/fhirfactory/harmonia/blob/98410db60155198a75f74c79d98060070209faa1/docs/markdown/05-information-systems-architecture/traceability/step-b-method.md\#b2-r07}{B2-R07} distinguishes individual clinical /\allowbreak{} health /\allowbreak{} care information from operational /\allowbreak{} workforce /\allowbreak{} financial /\allowbreak{} corporate /\allowbreak{} service-resource information. No exhaustive edge-case taxonomy or allocation to other enabling capabilities is introduced.

\phantomsection\label{rad:b2-1:138}

\texttt{\seqsplit{enabled-by}} does not mean that Health Information Access interprets, judges, plans, authorises or performs clinical care. It expresses material enablement, not a mandatory runtime dependency. B2.1 PASS does not change retained candidate deferrals, declare B1 PASS or authorise subsequent Step B mappings.
